(1) 轴对称的 $l=1$ Laplace 外解为 $\Phi=-E_0r\cos\theta+A\cos\theta/r^2$，不允许在无外加净电荷情况下另加 $1/r$ 单极项。球面切向电场零要求 $\Phi(a,\theta)$ 与角度无关，故 $A=E_0a^3$，球外势为 $\Phi=-E_0(r-a^3/r^2)\cos\theta$，可取球面电势为零。
(2) 径向场 $E_r=-\partial_r\Phi=E_0(1+2a^3/r^3)\cos\theta$，切向场 $E_\theta=E_0(1-a^3/r^3)\sin\theta$。在 $r=a^+$ 有 $E_r=3E_0\cos\theta$，故 $\sigma=3\varepsilon_0E_0\cos\theta$；对球面积分 $\int_0^\pi\cos\theta\sin\theta\,d\theta=0$，净荷严格为零。
(3) 真空点偶极势 $\Phi_{\rm dip}=p\cos\theta/(4\pi\varepsilon_0r^2)$；比较系数得 $p=4\pi\varepsilon_0a^3E_0$，静态极化率 $\alpha_e=p/E_0=4\pi\varepsilon_0a^3$。偶极电势为 $r^{-2}$ 而偶极场为 $r^{-3}$，不能混淆两者衰减阶。
(4) 诱导偶极不是固定值，而是 $p(E)=\alpha_eE$，因此可逆建立外场所对应的极化有效能变化为 $U_{\rm pol}=-\int_0^{E_0}p(E)\,dE=-\alpha_eE_0^2/2$。若直接使用永久偶极公式 $-pE_0$ 会多计一倍，忽视建立偶极所需的电源功。
